Each of the following produced a wrong number during this study before it was
caught. We report them as technical content for anyone benchmarking attention
variants, with the measured magnitude of each error, because a benchmarking
result is only as trustworthy as the harness that produced it.

\paragraph{P1. A diagnostic \texttt{.item()} inside the timed region.} The
module computed a mean-attention-entropy statistic in its forward pass. The
call synchronises the device \emph{inside} the region under measurement, so the
benchmark measures the synchronisation, not the module. Symptom: dense
attention latency appeared \emph{flat in $N$} --- $0.20$\,ms at $N{=}6$ rising
only to $0.48$\,ms at $N{=}192$ --- which we initially and wrongly read as
evidence that the $\mathcal{O}(N^2)$ cost does not manifest. Fix: gate all
diagnostics behind a flag that is off while timing.

\paragraph{P2. Timing arms in sequential blocks on a shared device.} Running
all repetitions of arm $A$ and then all of arm $B$ makes the comparison hostage
to whatever else used the device during each block. Symptom: swings above
$100\%$ between identical configurations of the same script, and in one case a
$3.7\times$ discrepancy for the same arm at the same $N$. Fix: interleave arms
repetition by repetition so interference lands on all of them simultaneously,
and report ratios, which are paired, alongside absolute latency, which is not.
This reduced observed spread from $>100\%$ to typically $<2\%$.

\paragraph{P3. \texttt{torch.cuda.synchronize()} without a device argument.}
The no-argument form synchronises the \emph{current} device. Timing work
submitted to a non-default device therefore waits on the wrong queue. Symptom:
every kernel, at every size, reports the CUDA-event resolution floor of
$\approx0.4\,\mu$s --- a result that looks like a spectacular speedup rather
than an obvious error. Fix: pin the device and pass it explicitly to every
synchronisation.

\paragraph{P4. Benchmarking a masked implementation as though it were sparse.}
The distinction of Section~\ref{sec:three}. Symptom: the ``sparse'' arm is
uniformly slower than dense, which is correct for that implementation but says
nothing about the method, since \Mask{} executes every multiply--accumulate
dense does plus selection on top. A FLOPs counter based on $1 - k/(N-1)$ will
happily report a $91\%$ saving for code whose executed operation count is
identical to dense. Fix: assert that executed FLOPs differ between the two
paths, and that the sparse path issues no $(B,N,N)\times(B,N,d)$ product. We
pin both as regression tests.

\paragraph{P5. Comparing against an unfused dense baseline.} If the deployed
alternative is a fused kernel, an unfused reference flatters the sparse method
--- by $4.6\times$ at $N{=}512$ in our setting. Note also that the fused path
is invisible to standard FLOP counters, which report zero for it, so a
FLOPs-only comparison cannot see the baseline that matters. Fix: report both,
and check which SDPA backend actually serves the shape and dtype in question
(in \texttt{fp32} only the memory-efficient kernel was available to us;
FlashAttention requires half precision).

\paragraph{A general form.} P1 and P3 are instrumentation errors that inflate
or flatten the measurement; P2 is a design error that lets a shared environment
leak into the comparison; P4 and P5 are specification errors about \emph{what
is being compared to what}. The last two are the more dangerous, because the
resulting numbers are internally consistent and reproducible --- they are simply
answers to a different question than the one the paper asks. The defence is to
state the operation count each arm executes alongside its latency, which makes a
mismatch between the claimed and the executed algorithm visible immediately.
